\begin{mainquestion}{ceilinglight}{Household Circuit}
At my house, the electrical wiring for the room lighting uses 14-gauge solid copper wire (see \href{https://www.engineeringtoolbox.com/awg-wire-gauge-d_731.html}{wire gauge chart}). There is a light switch which controls an overhead light in the living room; the wiring goes straight up the wall (behind the sheetrock) \qty{1}{\meter} to the ceiling where it makes a right-angle bend and travels \qty{4}{\meter} to the center of the room where the light fixture is mounted.

% =========================================================
\begin{partquestion}{electrontime}
{Calculate the expected transit time of an electron from the switch to the light fixture at the rated maximum current of \qty{15}{\ampere}. Assume the density of mobile electrons in copper is $n = \qty{8.5e22}{\per\centi\meter\cubed}$.}

\begin{answer}
Using the wire gauge chart link, we can see that that the 14-gauge wire is \qty{0.163}{\centi\meter} in diameter, which corresponds to a cross-sectional area A: 
\[
A = \frac{\pi d^2}{4}
  = \frac{\pi (\qty{0.163}{\centi\meter})^2}{4}
  = \qty{0.0209}{\centi\meter\squared}
\]

The formula for the average drift velocity (magnitude) for the electron is given by:
\[
\left|v_d\right|=\frac{I}{qnA}=\frac{15\mathrm{\ \operatorname{A}} }{\left(1.6\times10^{-19}\mathrm{\ C} \right)\left(8.5\times10^{22}\mathrm{\ c} \mathrm{\operatorname{m}}^{-3}\right)\left(0.0209\mathrm{\ c}\mathrm{\operatorname{m}}^2\right)}=0.0529\ \frac{\mathrm{cm} }{s}
\]

If the total length of wire is $\qty{1}{\meter} + \qty{4}{\meter} = \qty{5}{\meter} = \qty{500}{\centi\meter}$ , the transit time t is:
\[
t = \frac{\qty{500}{\centi\meter}}{\qty{0.0529}{\centi\meter\per\second}}
  = \qty{9.45e3}{\second}
\]

That's a pretty long time, about \qty{2.6}{\hour}!
\end{answer}
\end{partquestion}

% =========================================================
\begin{partquestion}{explaintime}
{Explain and account for any difference between your result in (\cref{prob:ceilinglight:electrontime}) and your real-world experience with how fast a light turns on after flipping the switch.}

\begin{answer}
Electrons interact through the Coulomb (repulsion) force. When one electron is disturbed, this creates a disturbance in the force (no kidding) which propagates as an electromagnetic wave. Electromagnetic waves propagate at the speed of light in vacuum and somewhat slower inside solid media. Ultimately, the speed of a signal along a wire is also limited by capacitive coupling with its environment, e.g., in a coaxial cable this is about 2/3 the speed of light in vacuum.

We also can think of a wire full of electrons like a pipe full of water. When we cause a current by pushing some water into one end of the pipe, water immediately starts coming out of the other end–we do not have to wait for the specific molecules we put in one end to pour out the other end. The speed of the water molecules does not matter so much as the speed with which an applied force is transmitted among the individual molecules. Both for water molecules and electrons, this force is transmitted so fast as to make the response appear nearly instantaneous.
\end{answer}
\end{partquestion}

% =========================================================
\begin{partquestion}{resistivity}
{Calculate the resistance of the copper wire that runs from the light switch to the overhead light. The resistivity of copper is $\rho = \qty{1.7e-8}{\ohm\meter}$. How much power is dissipated in the wire if the maximum \qty{15}{\ampere} is flowing?}

\begin{answer}
The resistance is given by
\[
R = \frac{\rho L}{A}
\]

Substitute values:
\[
R = \frac{(\qty{1.7e-8}{\ohm\meter})(\qty{5}{\meter})}{\qty{0.0209}{\centi\meter\squared}}
\]

Convert area to SI units:
\[
\qty{0.0209}{\centi\meter\squared} = \qty{2.09e-6}{\meter\squared}
\]

Now compute:
\[
R = \frac{(1.7 \times 10^{-8})(5)}{2.09 \times 10^{-6}}
  = \qty{0.0407}{\ohm}
  = \qty{40.7}{\milli\ohm}
\]

Power dissipation:
\[
P = I^2 R
\]

\[
P = (\qty{15}{\ampere})^2 (\qty{0.0407}{\ohm})
  = \qty{9.16}{\watt}
\]

\end{answer}
\end{partquestion}

% =========================================================
\begin{partquestion}{continuous}
{Typically, continuous loads in a housing circuit are limited to $80\%$ of the maximum current for safety. Find the power dissipated by the wire at this reduced current. How much power does this save compared to part (\cref{prob:ceilinglight:resistivity})?}

\begin{answer}
At $80\%$ of \qty{15}{\ampere}:
\[
I = 0.8 \times \qty{15}{\ampere} = \qty{12}{\ampere}
\]

Power:
\[
P_2 = I^2 R = (\qty{12}{\ampere})^2 (\qty{0.0407}{\ohm}) = \qty{5.86}{\watt}
\]

Power saved:
\[
\Delta P = \qty{9.16}{\watt} - \qty{5.86}{\watt} = \qty{3.30}{\watt}
\]

\end{answer}
\end{partquestion}

\end{mainquestion}